
%(BEGIN_QUESTION)
% Copyright 2015, Tony R. Kuphaldt, released under the Creative Commons Attribution License (v 1.0)
% This means you may do almost anything with this work of mine, so long as you give me proper credit

Les og lag en disposisjon av ``Introduction to Protective Relaying''''-delen i kapitlet ``Electric Power Measurement og Control'' i læreboken {\it Lessons In Industrial Instrumentation}. Noter sidetall der viktige illustrasjoner, fotografier, ligninger, tabeller og andre relevante detaljer finnes. Forbered deg på en faglig diskusjon med lærer og medstudenter om begrepene og eksemplene i lesingen.

\underbar{file i01247}
%(END_QUESTION)




%(BEGIN_ANSWER)


%(END_ANSWER)





%(BEGIN_NOTES)

Store effektbrytere må kommanderes til utløsning av eksterne enheter kalt {\it vernerelé}. ``Time-overcurrent''-reléer utløser brytere når målt strøm overstiger en grense, der utløsertiden avhenger av hvor mye grensen overskrides. ``Reclosing''-reléer kommanderer gjeninnkobling etter momentan overstrømsutløsning, for eksempel når en grein faller over en kraftlinje. Hensikten er å gjenopprette effekt hvis feilen var forbigående.

\vskip 10pt

125 VDC er standard forsyning for åpning/lukking av effektbrytere. En stor ``stasjonsbatteri''-bank opprettholder denne DC-forsyningen selv ved bortfall av nettkraft.

\vskip 10pt

Eldre vernerelé brukte induksjonsskive-elementer for å detektere utløsningsbetingelser og mekanisk lukke en kontakt som sender utløsersignal til bryterspole. Time-overcurrent-reléer er et godt eksempel. Slike reléer kunne ofte trekkes ut av sokkel for service og utskifting.

Senere elektroniske vernerelé brukte halvlederkomponenter for feildeteksjon og beslutning om utløsning. Mange av disse hadde også ``draw-out''-utførelse for enkel service.

Moderne mikroprosessorbaserte vernerelé er vanligvis panelmontert i stedet for draw-out, fordi behovet for regelmessig vedlikehold og utskifting er mindre. Disse reléene kan lagre data for senere analyse og kommunisere digitalt med andre systemer.

\vskip 10pt

Moderne digitale vernerelé bruker fortsatt enkelte historiske begreper som stammer fra elektrome-kaniske reléer.












\vskip 20pt \vbox{\hrule \hbox{\strut \vrule{} {\bf Forslag til sokratisk diskusjon} \vrule} \hrule}

\begin{itemize}
\item{} Beskriv funksjonen til et {\it vernerelé} med enkle ord.
\item{} Hva gjør et {\it reclosing}-relé, og hvorfor er denne funksjonen viktig i kraftdistribusjon?
\item{} Forklar why {\it reclosing} releer are only used for protecting overhead distribution lines, og never for protecting underground distribution lines.
\item{} Forklar hvorfor vernerelé vanligvis drives av 125 V {\it DC}. Hvorfor ikke 120 VAC?
\item{} Forklar what a {\it station battery} is in a substation or other electrical facility.
\item{} Beskriv hva ``pickup'' betyr i et vernerelé.
\item{} Hva betyr det at et vernerelé er av typen ``draw-out''?
\item{} Identifiser noen fordeler digitale vernerelé har sammenlignet med elektrome-kaniske reléer.
\item{} Er tidsforsinkelsen i et ``time-overcurrent''-relé fast, eller avhenger den av strømstyrken?
\item{} Studer det enkle relé/bryter-diagrammet i denne delen av læreboken, og identifiser flere mulige feilmodi (f.eks. komponent feilet kortsluttet eller åpen). Beskriv hvordan vernsystemet vil oppføre seg for hver feil.
\end{itemize}









\vfil \eject

\noindent
{\bf Forberedende quiz:}

Vis YouTube-videoen av kraftlinje som berører et tre i Bellingham, WA. Be deretter elevene forklare skriftlig hvorfor lysbuen {\it kom tilbake} etter første bryterutløsning.


%INDEX% Leseoppgave: Lessons In Industrial Instrumentation, introduksjon til vernerelésystemer

%(END_NOTES)

